\chapter{Introduction}

%% NOTE: English draft. main.tex now loads the class in `english' mode, so auto-generated
%% labels (Chapter, Figure, Table, References) render in English; the INTISARI page is
%% unaffected because \chapterintisari hardcodes its heading. Citation keys below
%% (e.g. \cite{virani2023}) are defined in references.bib.

\section{Background}

%% Scope note: population-scale burden is not reported at the CCD level -- CCD is a
%% clinical-management category introduced by \cite{virani2023}, not a Global Burden of
%% Disease cause. IHD, of which CCD is the chronic phase, is the closest reportable
%% level, so the figures below are IHD-specific rather than CVD-wide.
Chronic coronary disease (CCD), previously referred to as stable ischemic heart disease
(SIHD), is the chronic phase of ischemic heart disease (IHD), which carries the largest
single share of the global cardiovascular burden. The Global Burden of Disease Study
2019 estimated 197~million prevalent cases of IHD and 9.14~million IHD deaths worldwide
in 2019---out of 18.6~million cardiovascular deaths in total---and that burden continues
to rise in almost every country outside the high-income group \cite{gbd2019_cvd}.
Indonesia follows this pattern: between 2000 and 2019 its age-standardised cardiovascular
mortality rate rose from 356.05 to 412.46 deaths per 100{,}000 among men, while remaining
essentially unchanged among women \cite{gbd_indonesia_cvd}.
Patients with established CCD require lifelong secondary prevention,
and guideline-directed medical therapy (GDMT) has been shown to substantially reduce
the risk of major adverse cardiovascular events (MACE). To standardise this care, the American Heart Association and American College of Cardiology
published the 2023 AHA/ACC Guideline for the Management of Patients with Chronic Coronary
Disease \cite{virani2023}. It consolidates the evidence into a large set of recommendations,
spanning antiplatelet therapy, lipid and blood-pressure management, beta-blockers,
renin--angiotensin--aldosterone inhibition, weight and lifestyle management, and the safety
contraindications of those therapies.

Despite the availability of such guidelines, adherence in routine practice remains
suboptimal. The INTERASPIRE study, which examined 4{,}548 patients with coronary heart
disease across 14 countries in all six WHO regions, found that only 1.0\%
achieved its definition of optimal guideline adherence; just 38.6\% reached a
blood pressure below 130/80~mmHg, 16.6\% reached an LDL-cholesterol target of
below 1.4~mmol/L, and 9.0\% attended cardiac rehabilitation
\cite{adherence_gap}. Clinicians must reconcile a fragmented electronic
health record (EHR), dozens of coexisting conditions and medications, and a guideline
containing many conditional, patient-specific recommendations, often under severe time
pressure. Clinical decision support systems (CDSS) aim to close this gap by matching a
patient's record against the guideline and surfacing the applicable recommendations.
However, an effective CDSS for CCD must overcome three persistent challenges. First, it
must \emph{integrate} three distinct evidence sources---the patient's EHR, the clinical
guideline, and the supporting biomedical literature---into a single queryable
structure. Second, it must \emph{bridge heterogeneous clinical terminologies}
(e.g.\ SNOMED~CT, RxNorm, LOINC) so that a coded record can be matched to
guideline concepts. Third, when it communicates in natural language, it must be
\emph{trustworthy}, i.e.\ free of fabricated or unsupported clinical claims.

Two technological developments make a new approach feasible. First, modern EHRs increasingly expose clinical data through the HL7 Fast Healthcare
Interoperability Resources (FHIR) standard \cite{fhir_r4}. FHIR is a structured,
resource-oriented format. In it, conditions, procedures, medications, and \emph{numeric
observations} (laboratory results and vital signs) are first-class, machine-readable data. A \emph{FHIR-native}
CDSS can therefore consume real structured clinical data directly, rather than relying
on flattened tabular exports, in which a multi-component observation such as a
blood-pressure panel is dissolved into independent rows that must be reassembled before
they can be reasoned over (Section~\ref{sec:flattening}). That reassembly is a cost rather
than an obstacle; Section~\ref{sec:csv-gating} measures how large a cost. Second, knowledge
graphs (KGs), realised on property-graph
databases such as Neo4j, are well suited to fusing the EHR, the guideline, and the
literature into one graph, with terminology heterogeneity resolved by mapping every
code to a shared concept space using the Unified Medical Language System (UMLS)
\cite{umls}.

A complementary opportunity---and risk---arises from large language models (LLMs).
LLMs can render structured recommendation facts into a fluent, clinician-facing
briefing, but they are prone to hallucination, which is unacceptable in a clinical
setting. A design in which \emph{retrieval and reasoning are deterministic} (performed
by graph queries) and the LLM is restricted to \emph{phrasing} the already-retrieved,
provenance-carrying facts offers both readability and verifiable faithfulness, in
contrast to conventional retrieval-augmented generation that retrieves approximate
text passages via vector similarity \cite{rag_lewis2020}.

This thesis develops a FHIR-native, knowledge-graph-based clinical recommendation
system (CRS) for CCD. Synthetic but realistic patient records are generated with Synthea
\cite{synthea} and exported as HL7 FHIR R4 bundles, loaded into a Neo4j knowledge
graph, and fused with the 2023 AHA/ACC CCD guideline and a PubMed literature corpus
seeded from the guideline's references, through a UMLS concept layer. The
FHIR-to-graph transformation is itself routine and available as off-the-shelf,
web-based tooling \cite{tenasol_viewer}; this
work's contributions lie in the layers built above it. The central one follows
directly from the FHIR-native design. Because numeric observations are loaded as graph
data, the system can gate observation-dependent recommendations on the patient's actual
measured values. Examples are a beta-blocker recommendation conditioned on
left-ventricular ejection fraction and a weight-management recommendation conditioned on
body-mass index. A condition-only, text-based pipeline cannot do this. What produces the gain
is loading numeric observation values as first-class data, and FHIR is the most direct way to
obtain them---not the only way. Section~\ref{sec:csv-gating} tests this by rebuilding the same
gate from a tabular export, and Section~\ref{sec:limitations} states the bound that result puts
on the claim.
The system additionally produces a grounded natural-language briefing whose factual
content is retrieved deterministically from the graph and whose citation-level
faithfulness is checked automatically. In this way the work investigates how a FHIR-native knowledge
graph can improve both the precision and the trustworthiness of guideline-concordant
recommendations for chronic coronary disease.

\section{Problem Statement}

Based on the background above, the problems this research sets out to address are as
follows:

\begin{enumerate}
    \item The evidence needed to support a clinical decision in chronic coronary disease
    is spread across three sources of different shape---the patient record, the clinical
    practice guideline, and the biomedical literature---and is not available in any single
    queryable structure. Heterogeneous clinical terminologies compound this: conditions
    are coded in SNOMED~CT, medications in RxNorm, and observations in LOINC, so a coded
    record cannot be matched directly against the concepts a guideline refers to.

    \item A decision-support pipeline built on conditions and text alone discards numeric
    observation values, even though several guideline recommendations are conditioned on a
    patient's measured values---left-ventricular ejection fraction or body-mass index, for
    example. Observation-dependent recommendations therefore cannot be gated on what was
    actually measured, and the resulting guideline concordance is less precise than it
    could be.

    \item Natural-language output generated by an LLM is prone to hallucination, and
    conventional retrieval-augmented generation grounds only on approximately retrieved
    passages. No design is available whose faithfulness can be verified against the source
    record---a property that is not negotiable in a clinical setting.
\end{enumerate}

\section{Research Objectives}

In line with the problem statement, the objectives of this research are:

\begin{enumerate}
    \item To design and implement a FHIR-native pipeline that loads HL7 FHIR R4 bundles
    into a Neo4j knowledge graph and then fuses it with the 2023 AHA/ACC CCD guideline and
    a literature corpus seeded from the guideline's references, through a UMLS-based
    concept bridge. The resulting system is then evaluated end-to-end through guideline
    concordance against an automated, CRS-independent reference oracle
    (Section~\ref{sec:oracle}), agreement with a
    manually labelled gold standard (Cohen's~$\kappa$), contraindication-safety
    detection, and a clinician review of recommendation appropriateness.
    \item To implement observation gates that evaluate observation-dependent
    recommendations against the patient's measured values, and to quantify the
    resulting change in guideline concordance through an ablation study (with and
    without observation gating).
    \item To implement a grounded-briefing pipeline in which recommendation facts are
    retrieved deterministically from the graph and phrased by an LLM, and to evaluate
    its faithfulness using an automated verification procedure.
\end{enumerate}

\section{Research Scope}

To keep the research focused, the following scope and limitations apply:

\begin{enumerate}
    \item \textbf{Clinical domain.} The system targets chronic coronary disease
    (stable ischemic heart disease) only, and encodes the 2023 AHA/ACC CCD guideline
    \cite{virani2023} as its recommendation base.
    \item \textbf{Data.} Patient data are \emph{synthetic}, generated by Synthea and
    exported as HL7 FHIR R4 bundles; no real patient data are used. The cohort
    comprises 5{,}490 patients (see Chapter~III), of which the CCD
    gate-evaluable subset is used for evaluation.
    \item \textbf{Terminologies.} Conditions and procedures are coded in SNOMED~CT,
    medications in RxNorm, and observations in LOINC. ICD-10-CM / ICD-11 codes are not
    present in the synthetic data; ingesting ICD-coded records in a real deployment
    would require an additional SNOMED$\leftrightarrow$ICD crosswalk.
    \item \textbf{Observations.} Only a defined allow-list of clinically relevant LOINC
    observations (e.g.\ ejection fraction, LDL, BMI, triglycerides, blood pressure,
    smoking status) is loaded and used for gating.
    \item \textbf{Role.} The system is a \emph{recommender} and informational tool
    that surfaces guideline-concordant options; it does not make clinical decisions and
    is not deployed in live patient care.
    \item \textbf{Evaluation.} Effectiveness is measured as \emph{guideline
    concordance} rather than real clinical outcomes. Certain guideline preconditions
    that cannot be confirmed from structured data (e.g.\ composite risk scores and some
    negation-defined clauses) are treated as documented limitations.
\end{enumerate}

\section{Research Significance}

The expected benefits of this research are as follows.

\begin{enumerate}
    \item \textbf{Academic.}
    \begin{itemize}
        \item Demonstrates a reproducible architecture for a FHIR-native,
        knowledge-graph-based CRS that fuses EHR, guideline, and literature evidence.
        \item Quantifies the contribution of structured observations (enabled by the
        FHIR-native EHR) to guideline-concordance precision.
        \item Introduces a deterministic-retrieval faithfulness framework for LLM
        output in a clinical setting, in which factual content is verifiable against
        its source.
    \end{itemize}
    \item \textbf{Practical.}
    \begin{itemize}
        \item Provides a template for building a guideline-adherent recommender system on
        standardised FHIR data, portable to other conditions and guidelines.
        \item Offers a grounded-briefing pattern for trustworthy clinical text
        generation that can inform, without replacing, clinician judgement.
    \end{itemize}
\end{enumerate}

\section{Thesis Structure}

The remainder of this thesis is organised as follows.

\noindent\textbf{Chapter~I (Introduction)} presents the background, problem statement,
objectives, scope, and significance of the research.

\noindent\textbf{Chapter~II (Literature Review and Theoretical Foundation)} reviews
related work on clinical decision support, guideline knowledge graphs, and
retrieval-augmented generation, and establishes the theoretical basis, including HL7
FHIR, property-graph databases, UMLS-based terminology mapping, and the clinical
concepts of the CCD guideline.

\noindent\textbf{Chapter~III (Research Methodology)} describes the research design: the
data source and cohort, the graph construction order, the design of the
CRS-independent reference oracle, and the evaluation protocol.

\noindent\textbf{Chapter~IV (Results and Discussion)} first summarises the realised
system---the FHIR-to-Neo4j loader, the UMLS concept bridge, the guideline and literature
layers, the observation gates, the recommendation query layer, and the
grounded-briefing pipeline---then reports and discusses the graph construction quality,
the observation-gate concordance ablation, the manual $\kappa$-validation, the
safety-detection evaluation, and the briefing faithfulness results.

\noindent\textbf{Chapter~V (Conclusion)} summarises the findings, states the
conclusions with respect to the research objectives, and outlines directions for future
work.
